\section{Introduction}
Diffusion models learn generative transformations by connecting structured data to a tractable noisy reference distribution \citep{ho2020}. For directional observations, the sphere supplies a natural geometry, and intrinsic diffusion methods already extend classical score-based generation to manifolds \citep{debortoli2022}. A quantum formulation can likewise exploit a physical phase space rather than treating directions as arbitrary strings of basis labels. A spin-$j$ system offers precisely such a representation: its coherent states are indexed by the sphere, while the finite operator algebra resolves a hierarchy of angular multipoles.

Quantum diffusion itself is established prior work. QuDDPM learns ensembles of quantum data using scrambling and measurement-based denoising \citep{zhang2024}; QGDM constructs non-unitary channel-based generation of mixed-state targets \citep{chen2026qgdm}. Quantum-noise-driven generative models include quantum and hybrid forward/backward formulations \citep{parigi2024}, and mixed-state QuDDPM uses quantum noise channels and measured reverse circuits \citep{kwun2024}. QD3PM addresses discrete classical joint distributions through quantum diffusion \citep{chen2026qd3pm}. These precedents rule out identifying the present construction as the first quantum, fully quantum, or classical-data quantum diffusion model.

Our focus is the combination of a spin phase-space representation with a multipole-resolved physical diffusion and an auditable generative reconstruction. The coherent-state formalism \citep{arecchi1972} and quantum dynamical semigroups \citep{lindblad1976} are standard, as are phase-space methods for dissipative spin systems \citep{merkel2021}. We therefore treat the heat-spectrum correspondence as a mathematical foundation used by the model, not as a newly discovered property of angular momentum. The contribution is the explicit construction, its controlled training diagnosis, and a record connecting claims to reproducible numerical evidence. We make no priority claim for their combination based on the targeted literature search alone.

Three questions organize the experiments. First, which classical structures survive coherent-state encoding and readout at a given $j$? Second, once a secondary mode is representable, does failure to generate it reflect the noisy starting state, the reverse architecture, the objective, or the training budget? Third, does the product-phase-space extension learn correlations competitively in a small two-spin example? The first two questions yield positive but narrowly scoped results. The third yields a negative pilot result that limits extrapolation to many-body generation.

Throughout, ``quantum'' describes the state and channel model. Training, density evolution, and readout sampling are performed on a classical computer. Each single-spin run fits one empirical mixed state and generates directions through its Husimi distribution. It does not learn a distribution over individually labelled quantum training states in the ensemble-learning sense of QuDDPM.
